\chapter{实数直线上的拓扑}\label{ch:b1:topology}

极限始终借用同一套几何语言：点
“接近”某\omterm{def:b1:logic:sets}{集合}，\omterm{def:b1:logic:sets}{集合}“不漏\omterm{def:b1:topology:closure}{边界}”，数列无法逃出的
\omterm{prop:b1:reals:intervals}{区间}。本章确立这套术语——
\omterm{def:b1:topology:open}{开集}与\omterm{def:b1:topology:closed}{闭集}、\omterm{def:b1:topology:closure}{内部}与\omterm{def:b1:topology:closure}{闭包}、\omterm{def:b1:topology:dense}{稠密}性——在实数
直线上确立这些概念，并以序列形式证明闭\omterm{prop:b1:reals:intervals}{区间}的紧致性。
同样的概念在赋范向量空间中属于二年级内容；
在 $\R$ 上它们触手可及，并立即用于
\cref{ch:b1:continuity}。

\section{开集、闭集}

\begin{definition}[邻域、开集]\label{def:b1:topology:open}
\omterm{def:b1:logic:sets}{集合} $V \subseteq \R$ 称为\emph{邻域}\index{邻域}，
即 $x \in \R$ 的邻域，若它包含\omterm{prop:b1:reals:intervals}{区间} $\intoo{x - r}{x + r}$，
对某个 $r > 0$ 成立。\omterm{def:b1:logic:sets}{集合} $U \subseteq \R$ 称为\emph{开集}\index{开集}，
若它是其每一点的邻域：
\[
\forall x \in U,\ \exists r > 0, \quad \intoo{x - r}{x + r}
\subseteq U .
\]
\end{definition}

\begin{example}\label{ex:b1:topology:open}
开\omterm{prop:b1:reals:intervals}{区间}都是\omterm{def:b1:topology:open}{开集}：对 $x \in \intoo{a}{b}$，取 $r = \min(x -
a,\, b - x) > 0$。半直线 $\intoo{a}{+\infty}$ 是\omterm{def:b1:topology:open}{开集}；$\R$ 与
$\emptyset$ 是\omterm{def:b1:topology:open}{开集}（后者空真成立）。$\intcc{0}{1}$ 并
\emph{不是}\omterm{def:b1:topology:open}{开集}：$0$ 周围的任何\omterm{prop:b1:reals:intervals}{区间}都不完全含于其中。
\end{example}

\begin{proposition}[开集的稳定性]\label{prop:b1:topology:openstable}
任意多个\omterm{def:b1:topology:open}{开集}之并是\omterm{def:b1:topology:open}{开集}；\emph{有限}个\omterm{def:b1:topology:open}{开集}之交
是\omterm{def:b1:topology:open}{开集}。无限交集未必是\omterm{def:b1:topology:open}{开集}：
$\bigcap_{n \geq 1} \intoo{-\frac 1n}{\frac 1n} = \{0\}$，不是\omterm{def:b1:topology:open}{开集}。
\end{proposition}

\begin{proof}
并集情形：若 $x \in \bigcup_i U_i$，则 $x \in U_{i_0}$（$i_0$ 为某个指标），
由 $U_{i_0}$ 给出的\omterm{prop:b1:reals:intervals}{区间}含于并集之中。有限
交集：若 $x \in U_1 \cap \dots \cap U_k$，取 $r = \min(r_1,
\dots, r_k) > 0$ 为各 $U_j$ 所给出的半径。对于
反例：$0$ 的任一\omterm{prop:b1:reals:intervals}{区间}都含有某个 $\frac{1}{n}$
（阿基米德），因而不在交集之中。
\end{proof}

\begin{example}[用显式半径验证开性]\label{ex:b1:topology:radii}
$U = \{x \in \R : x^2 > 2\}$ 是\omterm{def:b1:topology:open}{开集}吗？答案是肯定的，且其证据
可以直接写出：$U = \intoo{-\infty}{-\sqrt2} \cup
\intoo{\sqrt2}{+\infty}$，为两个开半直线之并，由
\cref{prop:b1:topology:openstable}可知其为\omterm{def:b1:topology:open}{开集}。另可
逐点论证：对 $x \in U$ 且 $x > \sqrt 2$，取 $r = x -
\sqrt2 > 0$：对每个 $y \in \intoo{x - r}{x + r}$ 有 $y >
\sqrt 2$，从而 $y^2 > 2$；左侧情形同理。
两种风格都重要——结构性的（由已知\omterm{def:b1:topology:open}{开集}经
并集与有限交集构造）更易推广，
$\varepsilon$ 型适用于看不出结构的情形；而
\cref{ch:b1:continuity}将再给出第三种，也是最有力的方法：
$U$ 是开\omterm{prop:b1:reals:intervals}{区间} $\intoo{2}{+\infty}$ 在
连续\omterm{def:b1:logic:map}{映射} $x \mapsto x^2$ 下的\omterm{def:b1:logic:map}{原像}。
\end{example}

\begin{definition}[闭集]\label{def:b1:topology:closed}
\omterm{def:b1:logic:sets}{集合} $F \subseteq \R$ 称为\emph{闭集}\index{闭集}，当其
余集 $\R \setminus F$ 为\omterm{def:b1:topology:open}{开集}。由德摩根律与
\cref{prop:b1:topology:openstable}可知：任意多个闭集的交集是闭集，
有限个闭集的并集是闭集。
\end{definition}

\begin{theorem}[闭集的序列刻画]\label{thm:b1:topology:seqclosed}
$F$ 为\omterm{def:b1:topology:closed}{闭集}，当且仅当：对任一数列 $(u_n)$，若其各项属于 $F$，
则当其收敛于某个 $\ell \in \R$ 时，极限 $\ell$ 属于
$F$。（“闭”$=$“对极限封闭”。）
\end{theorem}

\begin{proof}
($\Rightarrow$) 设 $F$ 为\omterm{def:b1:topology:closed}{闭集}，$u_n \in F$，$u_n \to \ell$，且
假设 $\ell \notin F$。余集是\omterm{def:b1:topology:open}{开集}，于是存在
$\intoo{\ell - r}{\ell + r}$ 与 $F$ 不相交。但收敛性使 $u_n$
当 $n$ 足够大时落在该\omterm{prop:b1:reals:intervals}{区间}内，与 $u_n \in F$ 矛盾。

($\Leftarrow$) 设 $F$ 不是\omterm{def:b1:topology:closed}{闭集}：其余集不是\omterm{def:b1:topology:open}{开集}，
故存在 $x \notin F$，使得任何\omterm{prop:b1:reals:intervals}{区间} $\intoo{x - r}{x + r}$ 都不含于
余集之中；取 $r = \frac{1}{n+1}$，选取 $u_n \in F$ 满足
$\abs{u_n - x} < \frac{1}{n+1}$。于是 $u_n \in F$，$u_n \to x \notin
F$，序列性质不成立。
\end{proof}

\begin{example}[序列判别法的两个方向]\label{ex:b1:topology:seqtest}
\emph{\omterm{def:b1:topology:closed}{闭集}：}$F = \Z \cup \bigl\{n + \frac1n : n \geq 2\bigr\}$。
设 $u_k \in F$ 且 $u_k \to \ell$。取\omterm{prop:b1:reals:intervals}{区间} $\intcc{\ell -
1}{\ell + 1}$ 内仅含 $F$ 中有限多个点（有限多个
整数，有限多个 $n + \frac1n$），而从某一项起
所有 $u_k$ 都落在其中：此时数列只取有限多个
值，且由于收敛，最终为常数（参见
\cref{exo:b1:topology:3}）：$\ell \in F$。故为\omterm{def:b1:topology:closed}{闭集}——尽管
$F$ 中含有距离为 $\frac1n$ 的点对，任意
接近。

\emph{非\omterm{def:b1:topology:closed}{闭集}：}$G = \bigl\{\frac1m + \frac1n : m, n \in
\N^*\bigr\}$。数列 $\frac1n + \frac1n \in G$ 趋于
$0$，但 $0 \notin G$（两正数之和）：
序列判别法失效，$G$ 不是\omterm{def:b1:topology:closed}{闭集}。有趣的是，每个
$\frac1m$ \emph{确实}属于 $\overline G \cap G$：事实上
$\frac1m = \frac{1}{m+1} + \frac{1}{m(m+1)} \in G$。最后的要点是：
要证明闭性，须\emph{一并}控制所有收敛
数列（通常借助局部有限性或
闭式论证）；要否定它，只需一个选得恰当的逃逸
数列即可——这种不对称使得否定方向
较为容易，本章的反例都具有
这种一行即可写出的形式。
\end{example}

\begin{example}\label{ex:b1:topology:closed}
线段 $\intcc{a}{b}$、半直线 $\intco{a}{+\infty}$、有限
集、$\Z$（收敛的整数数列最终为常数）
都是\omterm{def:b1:topology:closed}{闭集}。$\intoc{0}{1}$ 既非\omterm{def:b1:topology:open}{开集}（在 $1$ 处失效），亦非
\omterm{def:b1:topology:closed}{闭集}（$\frac 1n \to 0 \notin$ 该\omterm{def:b1:logic:sets}{集合}）：多数\omterm{def:b1:logic:sets}{集合}两者皆非。$\R$ 与
$\emptyset$ 既是\omterm{def:b1:topology:open}{开集}又是\omterm{def:b1:topology:closed}{闭集}——而且它们是
$\R$ 中仅有的这类子集（见\cref{exo:b1:topology:9}）。
\end{example}

\begin{example}[由无穷多块拼成的开集]\label{ex:b1:topology:rminusz}
$\R \setminus \Z = \bigcup_{n \in \Z} \intoo{n}{n+1}$：这是
无穷多个开\omterm{prop:b1:reals:intervals}{区间}的并，由
\cref{prop:b1:topology:openstable} 知其为\omterm{def:b1:topology:open}{开集}——故 $\Z$ 为\omterm{def:b1:topology:closed}{闭集}，
无需任何数列论证。注意稳定性规则中的分工：
\omterm{def:b1:topology:open}{开集}的\emph{并}可以任意多（
每一点只需其自身的证书，由包含它的那个\omterm{def:b1:logic:sets}{集合}提供），
而\emph{交}必须是有限的（证书须求交，
且无穷多个半径可以缩为零）。
\Cref{exo:b1:topology:10} 将说明这一
例子即一般形式：$\R$ 的每个开子集都是
可数个不交开\omterm{prop:b1:reals:intervals}{区间}的并。
\end{example}

\section{内部、闭包、稠密性}

\begin{definition}[内部、闭包、边界]\label{def:b1:topology:closure}
设 $A \subseteq \R$。
\begin{itemize}
  \item 称点 $x$ 为 $A$ 的\emph{内点}，当 $A$ 是
        $x$ 的\omterm{def:b1:topology:open}{邻域}时；其\emph{内部}
        $\mathring{A}$\index{内部}是全体内点构成的\omterm{def:b1:logic:sets}{集合}。
  \item 称点 $x$ 为 $A$ 的\emph{附着点}，若
        $x$ 的每个\omterm{def:b1:topology:open}{邻域}都与 $A$ 相交；其\emph{闭包}
        $\overline{A}$\index{闭包}是全体附着点构成的\omterm{def:b1:logic:sets}{集合}。
  \item 其\emph{边界}为 $\partial A = \overline A \setminus
        \mathring A$。\index{边界}
\end{itemize}
于是 $\mathring A \subseteq A \subseteq \overline A$。
\end{definition}

\begin{proposition}[主要性质]\label{prop:b1:topology:closureprops}
\begin{enumerate}
  \item $\mathring A$ 是包含于 $A$ 的最大\omterm{def:b1:topology:open}{开集}；$A$
        为\omterm{def:b1:topology:open}{开集}，当且仅当 $A = \mathring A$。
  \item $\overline A$ 是包含 $A$ 的最小\omterm{def:b1:topology:closed}{闭集}；$A$
        为\omterm{def:b1:topology:closed}{闭集}，当且仅当 $A = \overline A$。
  \item （附着点的数列刻画）$x \in \overline
        A$ 当且仅当 $x$ 是某个由
        $A$ 中的点组成的数列的极限。
  \item 取补运算将两个概念互换：$\R \setminus
        \overline A = \bigl(\R \setminus A\bigr)^{\!\circ}$。
\end{enumerate}
\end{proposition}

\begin{proof}
(4) $x \notin \overline A$ $\iff$ 存在 $x$ 的某个\omterm{def:b1:topology:open}{邻域}与 $A$ 不相交
$\iff$ 存在包含 $x$ 的\omterm{prop:b1:reals:intervals}{区间}落在 $\R \setminus A$ 中 $\iff$ $x$
为 $\R \setminus A$ 的内点。

(1) $\mathring A$ 为\omterm{def:b1:topology:open}{开集}：若 $x \in \mathring A$，则存在
$\intoo{x-r}{x+r} \subseteq A$；该\omterm{prop:b1:reals:intervals}{区间}中每点 $y$ 都有
一个包含于其中的更小\omterm{prop:b1:reals:intervals}{区间}（从而包含于 $A$），因此整个
\omterm{prop:b1:reals:intervals}{区间}都在 $\mathring A$ 中。任一\omterm{def:b1:topology:open}{开集} $U \subseteq A$ 中的点都是
$A$ 的内点，故 $U \subseteq \mathring A$：它是最大的。
\omterm{def:b1:topology:open}{开集}的刻画随之可得。

(2) 详细推导。由 (4)，$\R \setminus \overline A$ 是
$\R \setminus A$ 的\omterm{def:b1:topology:closure}{内部}，由 (1) 为\omterm{def:b1:topology:open}{开集}，故 $\overline A$ 为
\omterm{def:b1:topology:closed}{闭集}，且包含 $A$。最小性：设 $F \supseteq A$ 为
\omterm{def:b1:topology:closed}{闭集}。则 $\R \setminus F$ 是\omterm{def:b1:topology:open}{开集}，且包含于 $\R
\setminus A$ 中，由 (1) 中的最大性，
\[
\R \setminus F \subseteq \bigl(\R \setminus A\bigr)^{\!\circ}
= \R \setminus \overline A ,
\]
再取一次补，得 $\overline A \subseteq F$。于是
$\overline A$ 是最小的闭超集。刻画如下：
若 $A = \overline A$，则 $A$ 为\omterm{def:b1:topology:closed}{闭集}（刚已证）；若 $A$ 为
\omterm{def:b1:topology:closed}{闭集}，则它本身即为包含 $A$ 的闭超集，故由最小性
得 $\overline A \subseteq A$，从而相等。

(3) 若 $u_n \in A$，$u_n \to x$：$x$ 的每个\omterm{def:b1:topology:open}{邻域}都包含
某个 $u_n \in A$，故 $x \in \overline A$。反之，若 $x \in
\overline A$：对每个\omterm{prop:b1:reals:intervals}{区间} $\intoo{x - \frac{1}{n+1}}{x +
\frac{1}{n+1}}$ 都与 $A$ 相交于某个 $u_n$，且 $u_n \to x$。
\end{proof}

\begin{example}\label{ex:b1:topology:closureexamples}
$\overline{\intoo{0}{1}} = \intcc{0}{1}$;
$\mathring{\intcc{0}{1}} = \intoo{0}{1}$; $\partial\intoo{0}{1} =
\{0, 1\}$。对于 $A = \{\frac 1n : n \in \N^*\}$：$\overline A = A
\cup \{0\}$, $\mathring A = \emptyset$, $\partial A = A \cup \{0\}$.
对于 $\Q$：由\omterm{def:b1:topology:dense}{稠密}性（\cref{thm:b1:reals:density}）每个实数都是
$\Q$ 的附着点，故 $\overline{\Q} = \R$，而 $\mathring{\Q} =
\emptyset$（每个\omterm{prop:b1:reals:intervals}{区间}都含无理数）。$\Q$ 的\omterm{def:b1:topology:closure}{边界}
为整个 $\R$。
\end{example}

\begin{example}[完整剖析]\label{ex:b1:topology:anatomy}
设 $A = \intoc{0}{1} \,\cup\, \bigl(\Q \cap \intoo{2}{3}\bigr)
\,\cup\, \{4\}$。下面逐片计算
\cref{def:b1:topology:closure} 中的三个\omterm{def:b1:logic:sets}{集合}。

\emph{\omterm{def:b1:topology:closure}{内部}。}$\intoo{0}{1}$ 中的任一点在
$A$ 中都有一整个\omterm{prop:b1:reals:intervals}{区间}，故为内点。点 $1$：其周围每个\omterm{prop:b1:reals:intervals}{区间}都
伸向 $1$ 的右侧，而 $A$ 在 $2$ 之前一无所有，故此点不是内点。
$\Q \cap \intoo{2}{3}$ 中的点都不是内点（每个
\omterm{prop:b1:reals:intervals}{区间}都含无理数，见\cref{thm:b1:reals:density}）；
孤立点 $4$ 也不是。于是 $\mathring A = \intoo{0}{1}$。

\emph{\omterm{def:b1:topology:closure}{闭包}。}$A$ 中各点的极限有：$\intcc{0}{1}$ 中的全部点，
（$0 = \lim \frac1n$，其中 $\frac 1n \in A$）；
$\intcc{2}{3}$ 中的全部点（该\omterm{prop:b1:reals:intervals}{区间}中每个实数都是其中有理数的极限，
即\omterm{def:b1:topology:dense}{稠密}性再次起作用）；以及 $4$。除此之外别无他点：
位于 $\intcc{0}{1} \cup \intcc{2}{3} \cup \{4\}$ 之外的点，
到该\omterm{def:b1:topology:closed}{闭集}的距离为正。于是 $\overline A = \intcc{0}{1} \cup
\intcc{2}{3} \cup \{4\}$。

\emph{\omterm{def:b1:topology:closure}{边界}。}$\partial A = \overline A \setminus \mathring A
= \{0, 1\} \cup \intcc{2}{3} \cup \{4\}$。

小结：三种运算都是\emph{局部}的——
每一片 $A$ 都依其自身性质作出贡献（
实心\omterm{prop:b1:reals:intervals}{区间}保留其\omterm{def:b1:topology:closure}{内部}，\omterm{def:b1:topology:dense}{稠密}而疏松的一片则整个化为\omterm{def:b1:topology:closure}{边界}，
孤立点则纯属\omterm{def:b1:topology:closure}{边界}），
而 $A$ 的两行草图
便可在任何证明写出之前预知每个答案。
\end{example}

\begin{remark}[点集推理中的常见误区]
(i) \emph{“非开”不等于“闭”}：多数\omterm{def:b1:logic:sets}{集合}两者皆非
（如 $\intoc{0}{1}$），另有两个\omterm{def:b1:logic:sets}{集合}两者皆是
（$\emptyset$，$\R$）——开与闭并非互为反义，而是
借助取补互为对偶。(ii) \emph{\omterm{def:b1:topology:closure}{内部}与\omterm{def:b1:topology:closure}{闭包}不可交换}：
取 $A = \Q$，
\[
\overline{\mathring A} = \overline\emptyset = \emptyset
\qquad\text{而}\qquad
\bigl(\,\overline A\,\bigr)^{\!\circ} = \mathring \R = \R :
\]
两种迭代算子的差异达到了\omterm{def:b1:logic:sets}{集合}所能有的最大程度。(iii)
\emph{无穷多个\omterm{def:b1:topology:closed}{闭集}之并未必是\omterm{def:b1:topology:closed}{闭集}}：
$\bigcup_{n\geq1} \intcc{\frac1n}{1} = \intoc{0}{1}$——恰为
交集反例的镜像，见
\cref{prop:b1:topology:openstable}。(iv) \emph{\omterm{def:b1:topology:dense}{稠密}并不意味着大}：
$\Q$ \omterm{def:b1:topology:dense}{稠密}、可数、\omterm{def:b1:topology:closure}{内部}为空，其补集也\omterm{def:b1:topology:dense}{稠密}；\omterm{def:b1:topology:dense}{稠密}性
只说明“任意接近每一点”，而非
“几乎全部”。周末问题中的
\omterm{pb:b1:topology:1}{康托尔集}（\cref{pb:b1:topology:1}）恰好说明相反的情形，
即一个拓扑上很小、却不可数地大的\omterm{def:b1:logic:sets}{集合}。
\end{remark}

\begin{example}[精确计算一个闭包]\label{ex:b1:topology:closurecomputed}
设 $G = \bigl\{\frac1m + \frac1n : m, n \in \N^*\bigr\}$（取自
\cref{ex:b1:topology:seqtest}）。断言：
\[
\overline G = G \,\cup\, \Bigl\{\frac1m : m \in \N^*\Bigr\}
\,\cup\, \{0\} .
\]
($\supseteq$) $\frac1m = \lim_n \bigl(\frac1m + \frac1n\bigr)$
且 $0 = \lim_n \frac2n$：由序列刻画可知皆为附着点。
（$\subseteq$）设 $x = \lim_k \bigl(
\frac{1}{m_k} + \frac{1}{n_k}\bigr)$；将每对排列，使
$m_k \leq n_k$。若 $(m_k)$ 无界，则某\omterm{def:b1:seq:subsequence}{子列}有 $m_k
\to \infty$，从而 $n_k \to \infty$ 亦然，且 $x = 0$。否则
$(m_k)$ 只取有限多个值，其中某一个（记为 $m$）
无穷多次出现；沿此\omterm{def:b1:seq:subsequence}{子列}，$\frac{1}{n_k} \to x -
\frac1m$：若 $(n_k)$ 有界，则它取某个值 $n$
无穷多次，于是 $x = \frac1m + \frac1n \in G$；若不然，则 $x =
\frac1m$。每种情形都落入所宣布的\omterm{def:b1:logic:sets}{集合}。本例要点：
求\omterm{def:b1:topology:closure}{闭包}本质上是一种紧性式的下标分情形分析：
下标有界意味着只取有限多个值（鸽巢原理），
下标无界则意味着极限逃逸——
答案呈现出极限点的典型两层结构：\omterm{def:b1:logic:sets}{集合}本身、其第一代极限，
以及它们的极限，
即 $0$。
\end{example}

\begin{definition}[稠密性，拓扑形式]\label{def:b1:topology:dense}
称 $A$ 在 $\R$ 中\emph{稠密}\index{稠密}，当 $\overline A = \R$
——也即每个非空开\omterm{prop:b1:reals:intervals}{区间}都与 $A$ 相交；
等价地（由\cref{prop:b1:topology:closureprops}~(3)），每个实数都是
$A$ 中元素的极限。例如 $\Q$，$\R \setminus
\Q$，二进有理数（见\cref{exo:b1:reals:8}），稠密\omterm{def:b1:structures:subgroup}{子群}
（见\cref{exo:b1:reals:9}）。
\end{definition}

\begin{example}[稠密性是相对的]\label{ex:b1:topology:relativedensity}
“\omterm{def:b1:topology:dense}{稠密}”按此处的定义，是指\emph{在 $\R$ 中}\omterm{def:b1:topology:dense}{稠密}；一个\omterm{def:b1:logic:sets}{集合}可以
反过来只在直线的一部分上\omterm{def:b1:topology:dense}{稠密}。$\intcc{0}{1}$ 的二进有理数，
即 $D \cap \intcc{0}{1}$，
（见\cref{exo:b1:reals:8}）与每个包含于
$\intcc{0}{1}$ 的开\omterm{prop:b1:reals:intervals}{区间}相交，但当然完全错过 $\intoo{2}{3}$：它们
\emph{在} $\intcc{0}{1}$ 中\omterm{def:b1:topology:dense}{稠密}，即 $\overline{D \cap
\intcc{0}{1}} = \intcc{0}{1}$。一般说法“$A$ 在 $B$ 中\omterm{def:b1:topology:dense}{稠密}”
是 $B \subseteq \overline A$ 的简写——务必指明
所在的背景\omterm{def:b1:logic:sets}{集合}，因为周末问题的端点是
\emph{在\omterm{pb:b1:topology:1}{康托尔集}中}\omterm{def:b1:topology:dense}{稠密}，同时却\emph{在
$\R$ 中}无处\omterm{def:b1:topology:dense}{稠密}：同一\omterm{def:b1:logic:sets}{集合}，两种皆真而听来相反的
描述。
\end{example}

\begin{example}[稠密性的处理]\label{ex:b1:topology:densitytricks}
三个反复出现的小技巧。\emph{放大}：若
$A$ \omterm{def:b1:topology:dense}{稠密}且 $A \subseteq B$，则 $B$ \omterm{def:b1:topology:dense}{稠密}（每个
\omterm{prop:b1:reals:intervals}{区间}本就与 $A$ 相交）。\emph{仿射变换}：若 $A$
\omterm{def:b1:topology:dense}{稠密}，则 $\lambda A + \mu$ 在 $\lambda \neq 0$ 时亦\omterm{def:b1:topology:dense}{稠密}——任一
\omterm{prop:b1:reals:intervals}{区间} $I$ 与 $\lambda A + \mu$ 相交，当且仅当\omterm{prop:b1:reals:intervals}{区间}
$\frac{I - \mu}{\lambda}$ 与 $A$ 相交，故
比如 $10^{-9}$ 的奇数倍亦\omterm{def:b1:topology:dense}{稠密}。\emph{求交则不然}：两个\omterm{def:b1:topology:dense}{稠密}
\omterm{def:b1:logic:sets}{集合}可以互不相交（$\Q$ 与 $\R \setminus \Q$）：
\omterm{def:b1:topology:dense}{稠密}性在并运算与仿射\omterm{def:b1:logic:map}{映射}下保持，而在交运算下从不保持。
\end{example}

\section{闭区间的紧性}

\begin{theorem}[闭区间是列紧的]\label{thm:b1:topology:compact}
设 $a \leq b$。$\intcc{a}{b}$ 中的任一点列都有
一个\omterm{def:b1:seq:subsequence}{子列}收敛到\emph{属于}
$\intcc{a}{b}$ 中的点。\index{列紧性}

更一般地，具有此性质（每个点列都有在该\omterm{def:b1:logic:sets}{集合}中收敛的\omterm{def:b1:seq:subsequence}{子列}）的 $\R$ 的
子集恰为
\emph{有界\omterm{def:b1:topology:closed}{闭集}}。
\end{theorem}

\begin{proof}
$\intcc{a}{b}$ 中的点列有界，故由波尔查诺--魏尔斯特拉斯
（\cref{thm:b1:seq:bw}）可取出一个收敛\omterm{def:b1:seq:subsequence}{子列}；其极限
仍属于 $\intcc{a}{b}$，因为闭\omterm{prop:b1:reals:intervals}{区间}是\omterm{def:b1:topology:closed}{闭集}
（\cref{thm:b1:topology:seqclosed}）。

一般情形。\emph{（有界\omterm{def:b1:topology:closed}{闭集} $\Rightarrow$ 紧）}：
设 $F$ 为有界\omterm{def:b1:topology:closed}{闭集}，$(u_n)$ 为 $F$ 中的点列。
$F$ 的有界性使点列有界，故由波尔查诺--魏尔斯特拉斯
可取出 $u_{\varphi(n)} \to \ell$；又 $\ell \in F$，因 $F$
为\omterm{def:b1:topology:closed}{闭集}，该\omterm{def:b1:seq:subsequence}{子列}是收敛点列
且各项属于 $F$（\cref{thm:b1:topology:seqclosed}）：两个假设各用一次，
有界性用于保证极限存在，
闭性用于保证极限属于该\omterm{def:b1:logic:sets}{集合}。\emph{（紧 $\Rightarrow$ 有界\omterm{def:b1:topology:closed}{闭
集}）}：若 $F$ 无界，取 $u_n \in F$ 使 $\abs{u_n} \geq
n$；每个\omterm{def:b1:seq:subsequence}{子列}都无界，从而发散
（\cref{prop:b1:seq:first}）：故完全不存在收敛\omterm{def:b1:seq:subsequence}{子列}。若 $F$
不闭，取 $u_n \in F$ 使 $u_n \to \ell \notin F$
（\cref{thm:b1:topology:seqclosed}）：每个\omterm{def:b1:seq:subsequence}{子列}都收敛到
$\ell \notin F$，故无\omterm{def:b1:seq:subsequence}{子列}\emph{在} $F$ 中收敛。
\end{proof}

\begin{example}[嵌套紧集]\label{ex:b1:topology:nested}
这是定理的第一个应用。设 $K_0 \supseteq K_1 \supseteq
K_2 \supseteq \dots$ 均为非空紧集（即有界\omterm{def:b1:topology:closed}{闭集}），且都是
$\R$ 的子集。则 $\bigcap_n K_n \neq \emptyset$。事实上，取 $x_n
\in K_n$（对每个 $n$），点列位于紧集 $K_0$ 中，
故有\omterm{def:b1:seq:subsequence}{子列} $x_{\varphi(n)}$ 收敛到某点 $x$
（\cref{thm:b1:topology:compact}）。对固定的 $m$，
$x_{\varphi(n)}$（满足 $\varphi(n) \geq m$ 的项）都落在\omterm{def:b1:topology:closed}{闭集}
$K_m$ 中，故极限 $x$ 属于 $K_m$
（\cref{thm:b1:topology:seqclosed}）；由 $m$ 的任意性，$x \in
\bigcap_n K_n$。另有一个有用的补充：若\omterm{def:b1:topology:open}{开集} $U$ 包含
$\bigcap_n K_n$，则 $U \supseteq K_n$ 对某个 $n$ 成立——若不然，
对点 $x_n \in K_n \setminus U$ 应用同样的论证；极限
$x$ 应属于 $\bigcap K_n \subseteq U$，但 $U$ 为\omterm{def:b1:topology:open}{开集}，迫使
$x_{\varphi(n)} \in U$ 最终成立，这是矛盾。两个结论
在缺少紧性时均不成立：$\bigcap_n \intoo{0}{\frac
1n} = \emptyset$ 与 $\bigcap_n \intco{n}{+\infty} = \emptyset$。
本段要点：紧性能把无穷多个非空性断言组成的链条，
转化为单个极限点——它是
在取无穷交时依然有效的工具，而
\cref{pb:b1:topology:1}将两次依赖于它。
\end{example}

\begin{omfigure}
\begin{tikzpicture}[scale=9.5]
  \node[left, font=\small] at (-0.02,0) {$C_0$};
  \fill[omDef] (0,-0.011) rectangle (1,0.011);
  \node[left, font=\small] at (-0.02,-0.09) {$C_1$};
  \fill[omDef] (0,-0.101) rectangle (0.3333,-0.079);
  \fill[omDef] (0.6667,-0.101) rectangle (1,-0.079);
  \node[left, font=\small] at (-0.02,-0.18) {$C_2$};
  \fill[omDef] (0,-0.191) rectangle (0.1111,-0.169);
  \fill[omDef] (0.2222,-0.191) rectangle (0.3333,-0.169);
  \fill[omDef] (0.6667,-0.191) rectangle (0.7778,-0.169);
  \fill[omDef] (0.8889,-0.191) rectangle (1,-0.169);
  \node[left, font=\small] at (-0.02,-0.27) {$C_3$};
  \fill[omDef] (0,-0.281) rectangle (0.0370,-0.259);
  \fill[omDef] (0.0741,-0.281) rectangle (0.1111,-0.259);
  \fill[omDef] (0.2222,-0.281) rectangle (0.2593,-0.259);
  \fill[omDef] (0.2963,-0.281) rectangle (0.3333,-0.259);
  \fill[omDef] (0.6667,-0.281) rectangle (0.7037,-0.259);
  \fill[omDef] (0.7407,-0.281) rectangle (0.7778,-0.259);
  \fill[omDef] (0.8889,-0.281) rectangle (0.9259,-0.259);
  \fill[omDef] (0.9630,-0.281) rectangle (1,-0.259);
  \node[below, font=\small] at (0,-0.30) {$0$};
  \node[below, font=\small] at (0.3333,-0.30) {$\tfrac13$};
  \node[below, font=\small] at (0.6667,-0.30) {$\tfrac23$};
  \node[below, font=\small] at (1,-0.30) {$1$};
  \draw[omThm, thick, -Stealth] (0.5,-0.042) -- (0.5,-0.062);
  \draw[omThm, thick, -Stealth] (0.5,-0.132) -- (0.5,-0.152);
  \draw[omThm, thick, -Stealth] (0.5,-0.222) -- (0.5,-0.242);
\end{tikzpicture}

{\small 中间三分之一构造的前四个阶段：
$C_n$ 的每个线段都去掉其开的中间三分之一，留下
$2^{n+1}$ 个线段构成 $C_{n+1}$，其总长为
$(\frac23)^{n+1}$。\omterm{pb:b1:topology:1}{康托尔集} $C = \bigcap_n C_n$——
\cref{pb:b1:topology:1} 所研究的对象——
正是上述嵌套紧集论证所保证的非空紧“残余”：
长度为零，却仍有不可数多个点
留存下来。}
\end{omfigure}

\begin{remark}
这是极值定理背后的引擎
（\cref{ch:b1:continuity}）以及海涅关于一致连续性的定理。
“紧”一词将获得其一般的（基于覆盖的）定义，
这要等到本科二年级。在 $\R$ 上，列紧性就足够了，
而“紧 $=$ 闭 $+$ 有界”是需要记住的结论。
\end{remark}

\begin{example}[并集下的边界]\label{ex:b1:topology:boundaryunion}
恒有 $\partial(A \cup B) \subseteq \partial A \cup \partial
B$：$\partial(A \cup B)$ 中的点，其每个\omterm{def:b1:topology:open}{邻域}
都与 $A \cup B$ 相交（从而与 $A$ 或 $B$ 相交，
且无穷多次与其中之一相交），并与 $A \cup B$ 的补集相交，该补集
同时位于两者的补集之中——经简短核对，即可把该点
归入 $\partial A$ 或 $\partial B$。这一包含可以是极为
严格的：取 $A = \Q$ 与 $B = \R\setminus\Q$，
\[
\partial(A \cup B) = \partial \R = \emptyset ,
\qquad
\partial A \cup \partial B = \R \cup \R = \R :
\]
两个参差不齐的\omterm{def:b1:logic:sets}{集合}可以粘合成一个无缝\omterm{def:b1:logic:sets}{集合}，它们的\omterm{def:b1:topology:closure}{边界}
彼此抵消。本段要点：\omterm{def:b1:topology:closure}{内部}与
\omterm{def:b1:topology:closure}{闭包}在并运算与交运算下表现得\emph{单调}，
但\omterm{def:b1:topology:closure}{边界}不然——应把 $\partial$ 看作一个导出量
（$\overline A \setminus \mathring A$），绝不视为
具有独立代数结构的算子。
\end{example}

\begin{remark}[本分册中的视角]
这里建立的词汇在本书中还将被用到两次。
在\cref{ch:b1:continuity}中，每条定理实际上都是拓扑
命题的伪装：介值定理说，
连续\omterm{def:b1:logic:map}{映射}保持\omterm{prop:b1:reals:intervals}{区间}性质；最值
定理说，连续\omterm{def:b1:logic:map}{映射}保持紧性，而证明
中直接按名称引用\cref{thm:b1:topology:seqclosed,thm:b1:topology:compact}。
在\cref{ch:b1:multivar}中，同样的定义被
移到 $\R^2$ 中重新审视，以圆盘代替\omterm{prop:b1:reals:intervals}{区间}：\omterm{def:b1:topology:open}{开集}、
\omterm{def:b1:topology:closure}{闭包}与紧性可逐字迁移，而
二元最值定理仍依赖于
波尔查诺--魏尔斯特拉斯（逐坐标提取\omterm{def:b1:seq:subsequence}{子列}）。唯一
\emph{不}能轻松推广的概念是
\omterm{prop:b1:reals:intervals}{区间}本身——在平面上，连通性取代
了凸性，这一故事始于
\cref{exo:b1:topology:9}中“只有 $\emptyset$ 与 $\R$ 既开
又闭”的结论。
\end{remark}

\section{练习}

\begin{exercise}[$\star$]\label{exo:b1:topology:1}
对下列每个\omterm{def:b1:logic:sets}{集合}，判断它是\omterm{def:b1:topology:open}{开集}、\omterm{def:b1:topology:closed}{闭集}、既开又闭，还是既非开又非闭
（并说明理由）：$\intoo{0}{1} \cup \intoo{2}{3}$；$\;\intco{0}{1}$；
$\;\{0\} \cup \intcc{1}{2}$; $\;\R \setminus \Z$; $\;\Q \cap
\intoo{0}{1}$.
\end{exercise}

\begin{exercise}[$\star$]\label{exo:b1:topology:2}
求下列\omterm{def:b1:logic:sets}{集合}的 $\mathring A$、$\overline A$ 与 $\partial A$：
$A = \intoc{0}{1} \cup \{2\}$; $\;A = \R \setminus \Q$;
$\;A = \bigl\{\frac{(-1)^n n}{n+1} : n \in \N\bigr\}$.
\end{exercise}

\begin{exercise}[$\star$]\label{exo:b1:topology:3}
证明有限集是\omterm{def:b1:topology:closed}{闭集}，先利用补集，
再借助数列刻画。
\end{exercise}

\begin{exercise}[$\star$]\label{exo:b1:topology:4}
证明：对任意 $A, B \subseteq \R$，有 $\overline{A \cup B} =
\overline A \cup \overline B$。举例说明
$\overline{A \cap B}$ 可能不同于 $\overline A \cap \overline
B$.
\end{exercise}

\begin{exercise}[$\star\star$]\label{exo:b1:topology:5}
设 $u_n \to \ell$（$\R$ 中），证明\omterm{def:b1:logic:sets}{集合} $\{u_n : n \in \N\}
\cup \{\ell\}$ 是\omterm{def:b1:topology:closed}{闭集}（若再补充有界性，便是紧集，
  有界——而它确实有界）。
\end{exercise}

\begin{exercise}[$\star\star$]\label{exo:b1:topology:6}
设 $U$ 为\omterm{def:b1:topology:open}{开集}，$A$ 为任意\omterm{def:b1:logic:sets}{集合}，证明 $U + A = \{u + a\}$
是\omterm{def:b1:topology:open}{开集}。由此推出，\omterm{def:b1:topology:open}{开集}与任一\omterm{def:b1:logic:sets}{集合}之和仍为\omterm{def:b1:topology:open}{开集}，
并作对比：举出两个\emph{\omterm{def:b1:topology:closed}{闭集}}，其和不是\omterm{def:b1:topology:closed}{闭集}。
\emph{（可试取 $\Z$ 与 $\sqrt 2\,\Z$，并结合\cref{exo:b1:reals:9}。）}
\end{exercise}

\begin{exercise}[$\star\star$]\label{exo:b1:topology:7}
称点 $x \in A$ 为 $A$ 的\emph{孤立点}，若
$x$ 的某个\omterm{def:b1:topology:open}{邻域}与 $A$ 仅交于 $x$。证明 $\Z$ 的每点都是
$\Z$ 中的孤立点，而 $A = \{\frac1n\}$ 的所有点皆孤立，但
$\overline A \neq A$，并证明各点皆孤立的\omterm{def:b1:logic:sets}{集合}
\omterm{def:b1:topology:closure}{内部}为空。
\end{exercise}

\begin{exercise}[$\star\star$]\label{exo:b1:topology:8}
证明有界集的\omterm{def:b1:topology:closure}{闭包}仍有界，并证明
$\sup \overline A = \sup A$（设 $A$ 非空且有\omterm{def:b1:reals:bounds}{上界}）。由此推出
$\sup A \in \overline A$：\omterm{def:b1:reals:bounds}{上确界}总是附着点。
\end{exercise}

\begin{exercise}[$\star\star\star$]\label{exo:b1:topology:9}
证明 $\R$ 中既开又闭的子集只有
$\emptyset$ 与 $\R$。\emph{提示：设 $A$ 既开又闭，且
$A \neq \emptyset$ 与 $\R \setminus A \neq \emptyset$；取 $a \in
A$、$b \notin A$，不妨设 $a < b$，考虑 $s = \sup\,(A \cap
\intcc{a}{b})$；判断 $s$ 能否属于 $A$ 或
其补集。}
\end{exercise}

\begin{exercise}[$\star\star\star$]\label{exo:b1:topology:10}
（\omterm{def:b1:topology:open}{开集}的结构）设 $U \subseteq \R$ 为非空\omterm{def:b1:topology:open}{开集}。对
$x \in U$，记 $I_x$ 为所有包含
$x$ 且含于 $U$ 的开\omterm{prop:b1:reals:intervals}{区间}之并。证明 $I_x$ 是开\omterm{prop:b1:reals:intervals}{区间}，
两个\omterm{def:b1:logic:sets}{集合} $I_x$、$I_y$ 相等或不交，并证明 $U$ 是
\emph{可数个}两两不交的开\omterm{prop:b1:reals:intervals}{区间}之并\emph{（
在每个\omterm{prop:b1:reals:intervals}{区间}中各取一个有理数）}。
\end{exercise}

\begin{exercise}[$\star\star$]\label{exo:b1:topology:11}
称点 $x \in \R$ 为 $A$ 的\emph{聚点}，若
$x$ 的每个\omterm{def:b1:topology:open}{邻域}都与 $A \setminus \{x\}$ 相交；这些点构成的\omterm{def:b1:logic:sets}{集合}
称为\emph{导集} $A'$。证明 $\overline A = A \cup A'$，
并证明 $A$ 为\omterm{def:b1:topology:closed}{闭集}当且仅当 $A' \subseteq A$。求
下列\omterm{def:b1:logic:sets}{集合}的 $A'$：$A = \{\frac 1n : n \in \N^*\}$、$A = \Z$ 与
$A = \Q$。
\end{exercise}

\begin{exercise}[$\star\star\star$]\label{exo:b1:topology:12}
设 $A \subseteq \R$ 为非空集，定义 $d_A(x) = \inf\{\abs{x - a}
: a \in A\}$。证明：
\begin{enumerate}
  \item $\abs{d_A(x) - d_A(y)} \leq \abs{x - y}$，对所有 $x, y$
        （$d_A$ 是 $1$-利普希茨 函数）；
  \item $d_A(x) = 0$ 当且仅当 $x \in \overline A$；特别地，
        若 $F$ 为\omterm{def:b1:topology:closed}{闭集}且 $x \notin F$，则
        $d_F(x) > 0$;
  \item 对每个 $\varepsilon > 0$，\omterm{def:b1:logic:sets}{集合} $V_\varepsilon =
        \{x : d_F(x) < \varepsilon\}$ 是\omterm{def:b1:topology:open}{开集}，包含 $F$，且
        $\bigcap_{\varepsilon > 0} V_\varepsilon = F$（$F$
        为\omterm{def:b1:topology:closed}{闭集}）：每个\omterm{def:b1:topology:closed}{闭集}都是可数个
        \omterm{def:b1:topology:open}{开集}的交。
\end{enumerate}
\end{exercise}

\section{问题：康托尔集，既小又庞大}

\begin{problem}[{周末问题——康托尔三分集：
长度为零、不可数、完备，且 $C + C =
\intcc{0}{2}$}]
\label{pb:b1:topology:1}
从 $\intcc{0}{1}$ 中去掉其开中间三分之一，再
去掉每个剩余线段的中间三分之一，并无限重复此过程：留下的
\omterm{def:b1:logic:sets}{集合}是\emph{\omterm{pb:b1:topology:1}{康托尔集}} $C$，即分析学中最基本的
反例工厂。本题构造它，
并借助 \cref{pb:b1:reals:1} 中的 $3$ 进制方法阅读它，
并刻画它悖论式的面貌：总长度为零，却
不可数；\omterm{def:b1:topology:closure}{内部}为空，却无孤立点；完全不连通，
但 $C + C$ 充满整个线段
$\intcc{0}{2}$。\index{康托尔集}形式地说：$C_0 = \intcc{0}{1}$，
        而 $C_{n+1}$ 由 $C_n$ 得到：删去 $C_n$ 的每个线段的开
        中间三分之一；最后 $C = \bigcap_{n \geq 0}
C_n$。全文中，\emph{三进制码}是指对 $x \in \intcc{0}{1}$ 而言的
        任一数字串 $(d_k)_{k\geq1}$，其中 $d_k \in \{0, 1, 2\}$，
        其值 $\sup_n \sum_{k=1}^n d_k 3^{-k}$ 等于 $x$——
        允许使用非正常码（最终均为 $2$）；由
        \cref{pb:b1:reals:1}（第 9--11 题）可知，每个 $x \in
\intcc{0}{1}$ 有一个或两个码，恰有两个码当且仅当 $x = m/3^N
\in \intoo{0}{1}$。

\medskip
\noindent\textbf{第一部分——构造。}
\begin{enumerate}
  \item 将 $C_1$ 与 $C_2$ 显式地写成若干线段的并，
        并用归纳法证明：$C_n$ 是
        $2^n$ 个长为 $3^{-n}$ 的闭线段的不交并。
  \item 证明 $C$ 是\omterm{def:b1:topology:closed}{闭集}、有界——从而为紧集
        （见\cref{thm:b1:topology:compact}）——非空，并证明
        $C_n$ 中每条线段的每个端点都属于
        $C$。
  \item $C_n$ 的总长度为 $\bigl(\frac23\bigr)^n$。
        由此推得：对每个 $\varepsilon > 0$，\omterm{def:b1:logic:sets}{集合} $C$ 可以
        被有限个总长度
        $\leq \varepsilon$ 的线段覆盖，即\omterm{pb:b1:topology:1}{康托尔集}的\emph{长度
        为零}。
  \item 证明 $C$ 所包含的任一\omterm{prop:b1:reals:intervals}{区间}长度 $\leq
        3^{-n}$ 对每个 $n$ 成立，因而该\omterm{prop:b1:reals:intervals}{区间}是单点集或空集：
        $\mathring C = \emptyset$。由于\omterm{def:b1:topology:closed}{闭集}且\omterm{def:b1:topology:closure}{内部}为空，
        $C$ 是\emph{无处\omterm{def:b1:topology:dense}{稠密}}集。
\end{enumerate}

\medskip
\noindent\textbf{第二部分——三进制码。}
\begin{enumerate}[resume]
  \item 证明自相似递推式
        \[
        C_{n+1} = \tfrac13 C_n \,\cup\,
        \bigl(\tfrac23 + \tfrac13 C_n\bigr),
        \qquad\text{因此}\qquad
        C = \tfrac13 C \,\cup\, \bigl(\tfrac23 + \tfrac13
        C\bigr),
        \]
        两部分互不相交：$C$ 由两份
        按比例 $\frac13$ 缩小的自身副本组成。
  \item 对 $n$ 用归纳法证明：$x \in C_n$ 当且仅当
        $x$ 有一个三进制码，其前 $n$ 位数字属于
        $\{0, 2\}$。由此推得：利用 $x$ 至多有
        两个码这一事实，$x \in C$ 当且仅当 $x$ 有一个
        \emph{不含数字 $1$}的码（称为\emph{无 $1$ 码}）。
  \item 码的应用：给出无 $1$ 码，分别对 $0$、$1$、
        $\frac13$、$\frac23$ 写出；证明 $\frac14 = (0.\overline{02})_3$
        且 $\frac34 = (0.\overline{20})_3$，故二者都属于
        $C$；并验证 $\frac14$ \emph{不是}
        任一 $C_n$ 的端点（端点形如 $m/3^n$）。
  \item 证明每个 $x \in C$ 恰有\emph{一个}无 $1$ 码
        \emph{（当 $x$ 有两个码时，证明两者中恰有一个
        含数字 $1$）}。结论：
        取值\omterm{def:b1:logic:map}{映射}是从 $\{0,2\}$ 串\omterm{def:b1:logic:sets}{集合}到
        $C$ 的\omterm{def:b1:logic:inj}{双射}。
  \item （对角线法）设 $k \mapsto x_k$ 为任一\omterm{def:b1:logic:map}{映射} $\N^* \to C$。
        构造一个 $\{0, 2\}$ 串，其第 $k$ 位
        与 $x_k$ 的码不同，由此推得 $C$ 不可数
——与之形成对比，第 3 题指出它
        在测度意义下可忽略。
\end{enumerate}

\medskip
\noindent\textbf{第三部分——拓扑面貌。}
\begin{enumerate}[resume]
  \item 综合目前的结论：$C$ 是紧集、不可数、
        长度为零、无处\omterm{def:b1:topology:dense}{稠密}。哪一个单一的包含关系
        $C \subseteq C_n$ 分别承载上述哪一条性质？
  \item （$C$ 是完备集）设 $x \in C$ 的无 $1$ 码为
        $(d_k)$。将数字 $d_n$（$0 \leftrightarrow
        2$）翻转，得到 $x_n \in C$，满足 $\abs{x_n - x} = 2 \cdot
        3^{-n}$。由此推得 $C$ 没有孤立点：每个
        $C$ 中的点都是 $C$ 中\emph{其他}点的极限。
  \item 证明第 2 题中的端点构成
        $C$ 的一个可数\omterm{def:b1:topology:dense}{稠密}子集\emph{（将编码截断至前 $n$
        位，之后以 $0$ 续写；可数性的证明与
        \cref{exo:b1:topology:10}相同）}。由此推得：
        $C$ 中典型的点——如 $\frac14$——\emph{不是}端点：
        端点是嵌在不可数的
        “躯体”之中的可数“骨架”。
  \item （完全不连通）设 $x < y$ 为 $C$ 中的两点。取 $n$
        使 $3^{-n} < y - x$，并构造点 $z \in
        \intoo{x}{y}$ 且 $z \notin C$。由此推得
        $C$ 中所含的非空\omterm{prop:b1:reals:intervals}{区间}只能是单点集。
\end{enumerate}

\medskip
\noindent\textbf{第四部分——$C$ 的算术性质。}
\begin{enumerate}[resume]
  \item 证明 $1 - C = C$ \emph{（\omterm{def:b1:logic:map}{映射} $x \mapsto 1 - x$ 作用于
        一个无 $1$ 码会怎样？回忆 $1 = (0.\overline{2})_3$）}。
  \item （码的加法）证明：若 $x$、$x'$ 分别有码
        $(a_k)$、$(b_k)$，则 $x + x' = \lim_n\,(t_n + t'_n)$，
        其中 $t_n, t'_n$ 为部分和。由此推得：每个 $y
        \in \intcc{0}{1}$ 是 $C$ 中两点的\emph{中点}。
        设 $(e_k)$ 为 $y$ 的一个码，选取数字
        $a_k, b_k \in \{0, 2\}$，使 $\frac{a_k + b_k}{2} =
        e_k$。
  \item 由此推得 $C + C = \intcc{0}{2}$，结合第 14 题，
        可得 $C - C = \intcc{-1}{1}$。具体地，将 $1$ 写成
        第 7 题中找到的两个非端点之和。
  \item 反思：一个长度为零的\omterm{def:b1:logic:sets}{集合}，其差集却充满
        $\intcc{-1}{1}$。“$C$ 在测度意义下可忽略”与
        “$C + C$ 占满全长”之间为何没有矛盾？
        （用一句话回答；想想长度控制的是什么，
        又不控制什么。）
\end{enumerate}

\medskip
\noindent\textbf{第五部分——有理元素与无理元素。}
\begin{enumerate}[resume]
  \item 结合第 8 题与
        \cref{pb:b1:reals:1}的周期性判别法：$C$ 中的一点是有理数，当且仅当
        其 不含数字 $1$ 的编码最终是周期的。
        用 $3$ 进制长除法验证 $\frac1{13} =
        (0.\overline{002})_3 \in C$。
  \item 给出 $C$ 中一个显式的\emph{无理}元素：
        取编码的值，其中 $d_k = 2$ 位于三角数位置
        $k = \frac{j(j+1)}{2}$，其余位置 $d_k = 0$。
        利用\cref{pb:b1:reals:1}中间隙增长的论证说明其无理性
        （第 20 题）。
  \item （\omterm{def:b1:logic:inj}{满射}到整段）考虑\omterm{def:b1:logic:map}{映射} $h$，它将
        $C$ 中 不含数字 $1$ 的编码为 $(d_k)$ 的点映到
        \emph{二进制}字符串 $\bigl(\frac{d_k}2\bigr)$ 的值，即\
        $h(x) = \sup_n \sum_{k=1}^n \frac{d_k}{2}\,2^{-k}$。
        证明 $h$ 把 $C$ \emph{\omterm{def:b1:logic:inj}{满射}}到 $\intcc{0}{1}$。
        由此，可忽略的 $C$ \omterm{def:b1:logic:inj}{满射}到一条全长线段上——
        这是 $C$ 不可数的第二种证明。
  \item （自相似长度）假设某种长度 $L$ 已为 $C$ 及其缩小副本定义，
        且与缩放相容，
        即 $L(\lambda A) = \lambda L(A)$、平移不变性，
        以及对第 5 题中不交分解的可加性。
        证明此时 $L(C) = \frac23\,L(C)$，
        迫使 $L(C) = 0$：仅凭自相似性便已
        判定 $C$ 的长度为零。
\end{enumerate}

\medskip
\noindent\textbf{第六部分——一个“胖”的近亲，以及寓意。}
\begin{enumerate}[resume]
  \item （“胖”\omterm{pb:b1:topology:1}{康托尔集}）重复构造，但只在第
        $n$ 步（$n = 0, 1, 2, \dots$）时，从当前的 $2^n$ 个
        线段各去掉一个长度为 $4^{-(n+1)}$ 的中央开\omterm{prop:b1:reals:intervals}{区间}。
        证明线段长度 $l_n$ 满足
        $l_{n+1} = \frac{l_n - 4^{-(n+1)}}{2}$，$l_n =
        \frac{2^n + 1}{2\cdot 4^n} > 0$，所得的
        $K = \bigcap K_n$ 为紧集且\omterm{def:b1:topology:closure}{内部}为空，去掉的总长度为
        $\sum_{n\geq0} 2^n
        4^{-(n+1)} = \frac12$。承认（直观的，属于本科三年级的）
        有限个\omterm{prop:b1:reals:intervals}{区间}并集的长度可加性，
        利用\cref{ex:b1:topology:nested}中的两条结论，
        证明覆盖 $K$ 的任意有限个开\omterm{prop:b1:reals:intervals}{区间}族的总长度
        $\geq \frac12$：$K$ 无处\omterm{def:b1:topology:dense}{稠密}，
        但\emph{不}可忽略。“小”有
        若干互不等价的含义。
  \item （距离）证明：对非空\omterm{def:b1:topology:closed}{闭集} $F \subseteq
        \R$ 与 $x \in \R$，\omterm{def:b1:reals:bounds}{下确界} $d(x, F)$
        \emph{可以取到}\emph{（极小化数列加波尔查诺--魏尔斯特拉斯定理）}。
        然后计算
        \[
        \max_{y \in \intcc{0}{1}} d(y, C) = \frac16 ,
        \]
        恰在中心 $y = \frac12$ 处取到 \emph{（在第 $n$ 步
        产生的缺口中的一点距该缺口的端点
        不超过 $\frac{3^{-n}}{2}$，而端点属于
        $C$）}。
  \item （每个点都是\omterm{def:b1:seq:subsequence}{子列}极限）利用第 12 题与第 2 题，
        构造 $C$ 中的一个数列，其\omterm{def:b1:seq:subsequence}{子列}极限的\omterm{def:b1:logic:sets}{集合}
        恰为\emph{全部} $C$。（比较：
        对收敛数列，该\omterm{def:b1:logic:sets}{集合}恰为一个点——
        在紧集中，$C$ 代表了相反的极端。）
  \item 综合归纳，每条用一句话回答：（i）本章的哪些定理
        在构造中实际用到了（\omterm{def:b1:topology:closed}{闭集}的稳定性、紧性、
        数列刻画）？（ii）列出概貌中的四组悖论性搭配
        （长度为零与不可数，
        \omterm{def:b1:topology:closed}{闭集}与\omterm{def:b1:topology:closure}{内部}为空，
        完全集与完全不连通，
        可忽略与 $C+C$ 为满长度）；（iii）$C$ 在后文
        何处再次出现（在连续性理论中，以 $h$ 构造的
        魔鬼阶梯，以及本科三年级分册的测度论，
        其中 $C$ 用来区分“可数”与
        “可忽略”）？
\end{enumerate}
\end{problem}
